The repository models progressive aortic-wall degradation driven by chronic gut-derived signals (TMAO/LPS → ROS/HIF-1α → MMP-mediated matrix breakdown) together with cyclic vibrational / turbulent loading that further fatigues the media and concentrates stress. This degradation process is formalized in the limit model of degradation_failure.py (and narratively in the accompanying .tex files), showing how remaining ultimate tensile strength can decline over years until local wall stress exceeds the reduced threshold.
Within that framework the biochemical arm (everyday dysbiosis → matrix weakening) remains associative with respect to clinical aortic dissection risk; the vibrational-stress / cavitation (or balloon-expansion / jet-hammer) arm remains a specialized biomechanical hypothesis for how a sufficiently degraded wall may ultimately tear. The models do not assert that these pathways operate routinely or spontaneously in the absence of established risk factors.
Established determinants—hypertension, medial degeneration, genetic aortopathies, bicuspid aortic valve, and related hemodynamic or structural abnormalities—continue to dominate epidemiology and pathogenesis. The repository’s contribution is therefore best read as an explanatory model of aortic degradation under vibrational stress and of the conditions under which a weakened wall may dissect, rather than as a claim of spontaneous, everyday initiation of human aortic dissection.
